\documentclass[11pt]{article}
\usepackage[margin=25mm]{geometry}
\usepackage{fontspec}
\setmainfont{Latin Modern Roman}
\usepackage{amsmath,amssymb,mathtools}
\usepackage{unicode-math}
\setmathfont{Latin Modern Math}
\usepackage{xurl,hyperref}
\hypersetup{colorlinks=true,urlcolor=blue}
\setcounter{secnumdepth}{0}
\setlength{\parindent}{0pt}
\setlength{\parskip}{5pt}
\setlength{\emergencystretch}{3em}
\newcommand{\sref}[2]{\hyperlink{source:#1:#2}{[S#2]}}
\newcommand{\cataloguescope}[1]{\paragraph{Recorded target.}#1}
\newcommand{\R}{\mathbb R}
\newcommand{\C}{\mathbb C}
\newcommand{\N}{\mathbb N}
\newcommand{\Z}{\mathbb Z}
\newcommand{\Q}{\mathbb Q}
\newcommand{\F}{\mathbb F}
\newcommand{\eps}{\varepsilon}
\begin{document}
% UnsolvedMath ID 30003582; source catalogue code OWR-15584-003
\setcounter{section}{30003581}
\section{Quantum Backflow Bounds with External Potentials}\label{30003582}
{\small UnsolvedMath ID 30003582 · Mathematical Physics}
\subsection{Definitions and mathematical statement}

\paragraph{Shared notation.}$\mathbb N=\{1,2,\ldots\}$; the standard number fields have their usual meanings. All additional conventions are stated locally.


\paragraph{Definitions and conventions.}
In one-dimensional quantum mechanics, the probability current of a wave function is $j_\psi(x)=\frac{\hbar}{m}\operatorname{Im}(\overline{\psi(x)}\partial_x\psi(x))$. A smeared flux integrates this current against a spatial test function or the source\textquotesingle s associated quadratic form. Positive incoming states are scattering states whose asymptotic incoming momentum is positive. Backflow is negative current despite that positive incoming momentum. A lower bound compares the flux quadratic form with the state norm uniformly over this positive-incoming subspace. The primary smearing convention is a nonnegative Schwartz function $f\in\mathcal S(\mathbb R)$ and $J_f=\int f(x)J(x)\,dx$ as a quadratic form. The scattering Hamiltonian is $H=P^2/2+V(X)$ in units $\hbar=m=1$.

\paragraph{Target retained from the canonical source catalogue.}
\paragraph{Catalogue wording retained for traceability.}
For a quantum particle scattering from an external potential, is the smeared probability-flux operator on positive incoming states bounded below, analogously to the free-particle quantum-backflow inequality?

\paragraph{Source scope and qualification.}
The potential class, smearing function and scattering-wave-operator assumptions are essential and must follow the report. Positivity of incoming momentum does not mean the interacting state has positive momentum support at every time.
\subsection{Short English statement}
Does quantum backflow remain bounded when a particle scatters through an external potential?
\subsection{Research attempt}
\paragraph{Model, date and investigated scope.}
OpenAI GPT-6 Astra (configured model \texttt{gpt-6-astra}), ultra reasoning; generated 11 October 2026. Derived an elementary free lower bound for square Schwartz smearings and isolated a weighted wave-operator estimate; the external-potential target remains unproved.

\paragraph{Sources and literature check.}
Read the supplied definition, catalogue wording, source qualifications and reference entries for 30003582. The historical source/context and inherited status are retained; no fresh literature or novelty certification is claimed. No external theorem is used without proof in this attempt; explicit model assumptions are distinguished from facts about the source model.

\paragraph{Free estimate and its scattering obstruction.} Work in the source units and restrict to smearing \(f=g^2\) with real Schwartz \(g\). Write \(P=-i\partial_x\) and use the unitary Fourier transform. For smooth \(\psi\), integration by parts gives the real flux
\[
 \langle J_{g^2}\rangle_\psi=\operatorname{Re}\langle\psi,g^2P\psi\rangle=\langle g\psi,Pg\psi\rangle.
\]
Assume \(\widehat\psi(q)=0\) for \(q<0\). Dropping the nonnegative momentum part and using the Hilbert--Schmidt bound for the remaining integral kernel gives
\[
 \langle J_{g^2}\rangle_\psi\ge-\frac{\|\psi\|_2^2}{2\pi}
 \int_{k<0,q>0}|k|\,|\widehat g(k-q)|^2\,dq\,dk
 =-\frac{\|\psi\|_2^2}{4\pi}\int_{s<0}s^2|\widehat g(s)|^2\,ds.
\]
Indeed set \(k=-a,q=b\) and \(u=a+b\); integrating \(a\) from zero to \(u\) gives \(u^2/2\). This is an elementary free backflow bound, not its sharp constant.
\paragraph{Attempted transfer.} For an incoming wave operator \(\Omega_-\), the interacting state is \(\Omega_-\psi\), and its Fourier support need not be positive. Repeating the proof would require boundedness of
\[
 |P|^{1/2}\mathbf1_{(-\infty,0)}(P)\,g\Omega_-\mathbf1_{(0,\infty)}(P).
\]
Unitarity or isometry of \(\Omega_-\) alone does not control this unbounded weighted operator. Establishing the bound from the source potential assumptions is the first actual scattering gap. The restriction \(f=g^2\) is stated because an arbitrary nonnegative Schwartz function need not have a Schwartz square root. No source potential class or interacting theorem is silently inferred, and the full nonnegative-smearing problem remains unanswered.

\paragraph{Verification and adversarial audit.}
Checked Fourier normalization, the triangular integral constant, the square-root restriction on smearings and the distinction between positive incoming and instantaneous momentum. This is a GPT-6 Astra Ultra self-audit, not independent human verification. Prompt A drafting followed by prompt B scope/dependency/boundary self-audit. Checked Fourier normalization, the triangular integral constant, the square-root restriction on smearings and the distinction between positive incoming and instantaneous momentum. The remaining gap is: Bound the negative-momentum weighted operator with the actual incoming wave operator under the source potential hypotheses. Unresolved status and zero original-target completion are retained.

Coordinator checked the unitary Fourier convention, Hilbert--Schmidt kernel factor 1/(2 pi), triangular integration factor 1/2, and the distinction between incoming and instantaneous momentum. The interacting weighted-operator bound remains unproved. This additional check was model-assisted, not independent human verification.

\paragraph{Tools, process and contributions.}
Python 3 and local file-reading/writing tools for manuscript serialization and structural checks; mathematical calculations were symbolic reasoning, with no numerical experiment or theorem prover. Contribution prompt A was followed by prompt B. GPT-6 Astra with ultra reasoning drafted the research and performed its stated self-audit; any additional coordinator contribution is identified in the verification and audit. A GPT-6-based Codex coordinator assembled the manuscript and checked repository integration; its exact variant and reasoning configuration were not exposed. No human mathematical contribution, independent human proof check or formal proof certification is supplied. Complete derivations and the consulted source references are included here; private reasoning transcripts are not published.

\paragraph{Final status and first remaining gap.}
\textbf{UNRESOLVED}. Bound the negative-momentum weighted operator with the actual incoming wave operator under the source potential hypotheses.

\section{Completion Estimate}
\noindent\textbf{COMPLETION: 0\%}
The proved calculation is an elementary auxiliary result or a test of a restricted strategy; no new portion of the full source target is established.

\subsection{Sources}
\begin{itemize}
\item[\hypertarget{source:30003582:1}{S1}] ULAM.AI and the UnsolvedMath Contributors, UnsolvedMath, record 30003582 (OWR-15584-003). Catalogue curation under CC BY 4.0; definitions and English exposition added locally.
\url{https://huggingface.co/datasets/ulamai/UnsolvedMath}.
\item[\hypertarget{source:30003582:2}{S2}] Mathematical Questions and Challenges in Quantum Electrodynamics and its Applications (2017), Oberwolfach Reports, DOI 10.4171/owr/2017/41, pp. 2539--2599. Primary report accessed 2026-10-11; relevant PDF pages 51, 52, 53.
\url{https://doi.org/10.4171/owr/2017/41}.
\item[\hypertarget{source:30003582:3}{S3}] Bostelmann--Cadamuro--Lechner (catalogue-listed context; not independently consulted in this preparation).
\url{https://doi.org/10.1103/PhysRevA.96.012112}.
\item[\hypertarget{source:30003582:4}{S4}] Cadamuro (catalogue-listed context; not independently consulted in this preparation).
\url{https://arxiv.org/abs/1710.04322}.
\end{itemize}
\label{end:30003582}
\end{document}
